\subsection*{Experimental items from Experiment 1:}
\begin{enumerate}
\item \textbf{context} Grupa przestępcza składała się wyłącznie z kobiet.\\
  \textbf{pre} Najkrótszy wyrok otrzymał/otrzymała szpieg.\\
  \textbf{post} Szpieg otrzymał/otrzymała najkrótszy wyrok.\\
  \textbf{pron} Szpieg, który/która pozostaje w areszcie, spodziewa się najkrótszego wyroku.\\
  Q: Czy w zdaniu powiedziano jaki dokładnie to był wyrok? A: Nie
\item Przy stajni kręci się kilka pracowniczek, jedna z nich na koniu.\\
	Spośród nich najlepiej lekcje jazdy konnej prowadziła jeździec.\\
	Q: Czy w zdaniu była mowa o lekcjach jazdy konnej? A: Tak
\item Lekarki nie są zgodne co do diagnozy.\\
	Na wpływ leków uspokajających na przebieg choroby uwagę zwróciła psychiatra.\\
	Q: Czy w zdaniu była mowa o lekach uspokajających? A: Tak
\item Lekarki otworzyły salę i zaprowadziły pacjenta do stołu operacyjnego.\\
	Ponieważ wcześniej było to niemożliwe, wenflon pacjentowi założyła chirurg.\\
	Q: Czy w zdaniu był wspomniany wenflon? A: Tak
\item Na wydarzeniu pojawiły się również reprezentantki gminy Sterdyń.\\
	Najlepsze wrażenie na zgromadzonych zrobiła wójt.\\
	Q: Czy zdanie opisało na czym polegało wydarzenie? A: Nie
\item Islandia jest jednym z niewielu państw na świecie, gdzie kobiety pełnią funkcje prezydenta i szefa rządu.\\
	W 2024 roku swoje stanowisko zajęła premier.\\
	Q: Czy premier jest na stanowisku od 2009 roku? A: Nie
\item Na coroczne Dni Ekologii zaprosiłyśmy pracowniczki Instytutu Biologii Ssaków PAN.\\
	Wiekszość wykładów była prowadzona przez nasze koło naukowe, ale wykład o żubrach i turystyce poprowadziła naukowiec.\\
	Q: Czy w zdaniu był wspomniany wykład na temat dzików? A: Nie
\item Ostatnio kilka kobiet w jednostce wojskowej awansowało na wyższe stanowiska.\\
	Obowiązki głównego księgowego objęła porucznik.\\
	Q: Czy w zdaniu był wspomniany awans? A: Tak
\item Tegoroczne praktykantki w zakładach chemicznych, z których wszystkie poza jedną mają tytuł doktora, powoli zaczynają swoje praktyki.\\
	W zeszłym tygodniu praktyki zaczęła magister.\\
	Q: Czy w zdaniu były wspomniane praktyki? A: Tak
\item Wszystkie siostry przyjechały razem na miejsce samochodem.\\
	Ogromny kosz piknikowy z bagażnika wyciągnęła kierowca.\\
	Q: Czy w zdaniu był wspomniany kosz piknikowy? A: Tak
\item Ceremonię otworzyły najbardziej szanowane pracowniczki wydziału.\\
	Przemową skierowaną do pierwszorocznych studentów zaczęła dziekan.\\
	Q: Czy przemowa była skierowana do studentów pierwszego roku? A: Tak
\item Ósmego marca wszystkie pracowniczki gminy otrzymały po bukiecie kwiatów.\\
	Największy bukiet otrzymała burmistrz.\\
	Q: Czy prezentem wspomnianym w zdaniu były pralinki? A: Nie
\item Grupa policjantek weszła do knajpy.\\
	Spośród nich najwięcej jedzenia zamówiła sierżant.\\
	Q: Czy w zdaniu wspomniano jakie jedzenie zostało zamówione? A: Nie
\item Kobiety wśród władz uczelni bardzo poważnie potraktowały skargę.\\
	Postępowanie dyscyplinarne wszczęła rektor.\\
	Q: Czy w zdaniu mowa o postępowaniu dyscyplinarnym? A: Tak
\item W jednostce wojskowej, w której znajdują się same kobiety, odbył się ostatnio apel z okazji Święta Wojska Polskiego.\\
	Z tej okazji uroczystą zbiórkę poprowadziła pułkownik.\\
	Q: Czy w zdaniu wspomniania była uroczysta zbiórka wojskowa? A: Tak
\item Wszystkie najważniejsze kobiety z administracji miasta pojawiły się na otwarciu nowej biblioteki.\\
	O znaczeniu tego miejsca dla zapobieganiu samotności osób starszych mówiła prezydent.\\
	Q: Czy w zdaniu były wspomniane osoby starsze? A: Tak
\item W tym tygodniu odwiedzają nas mamy uczniów z naszej klasy, żeby opowiedzieć o swojej pracy.\\
	We wtorek przyszła architekt.\\
	Q: Czy dzień tygodnia, który był wspomniany w zdaniu, to czwartek? A: Nie
\item Delegacja z Ministerstwa Kultury i Dziedzictwa Narodowego, która odwiedziła Panteon w Paryżu, składała się wyłącznie z kobiet.\\
	Na grobie Marii Skłodowskiej-Curie kwiaty złożyła minister.\\
	Q: Czy w zdaniu mowa o zapaleniu znicza na grobie? A: Nie
\item Nad projektem miasta pracował zespół składający się wyłącznie z kobiet.\\
	W biurze przy Rynku Głównym pracowała inżynier.\\
	Q: Czy w zdaniu był wspomniany Rynek Główny? A: Tak
\item Dopiero lekarki ze Szpitala Bielańskiego potraktowały sprawę poważnie.\\
	Skierowanie na laparoskopię wystawiła ginekolog.\\
	Q: Czy w zdaniu była wspomniana gastroskopia? A: Nie
\item Przy przystanku pod teatrem stała grupa kobiet, spośród których jedna była wcześniej na widowni.\\
	Kiedy tylko autobus przyjechał, do środka wsiadła widz.\\
	Q: Czy w zdaniu wspomniany był tramwaj? A: Nie
\item Grupa studentek przyszła z prośbą o przesunięcie terminu kolokwium.\\
	O odwołanych ćwiczeniach i konieczności odrobienia ich przed kolokwium przypomniała starosta.\\
	Q: Czy studentki w zdaniu chcą odwołać ćwiczenia? A: Nie
\item Na spotkanie przyszła grupa starszych kobiet, z których jedna brała udział w Powstaniu Warszawskim, żeby opowiedzieć o życiu w Polsce w czasie II wojny światowej. \\
	Na samym początku głos zabrała powstaniec.\\
	Q: Czy w zdaniu była mowa o osobie która brała udział w Powstaniu Warszawskim? A: Tak
\item Wczoraj po południu odbył się panel, w którym kobiety opowiadały o swoich doświadczeniach w zawodach zdominowanych przez mężczyzn.\\
	O przemocy ze względu na płeć opowiadała polityk.\\
	Q: Czy w zdaniu wspomniano o przemocy na tle rasowym? A: Nie
\end{enumerate}

\subsection*{Experimental items from Experiment 2:}
\begin{enumerate}
\item \textbf{pre} Na rowerach przyjechały/przyjechali chłopaki.\\
  \textbf{post} Chłopaki przyjechały/przyjechali na rowerach.\\
  \textbf{pron} Chłopaki, które/którzy jadą na rowerach, mają szansę przyjechać na czas.\\
	Q: Czy w zdaniu były wspomniane rowery? A: Tak
\item Wtedy gołębie i króliki chodowały górniki.\\
	Q: Czy w zdaniu była mowa o papugach? A: Nie
\item W tym hotelu mieszkały generały.\\
	Q: Czy w zdaniu była mowa o hotelu? A: Tak
\item Zakupu tych opon mi odradziły mechaniki.\\
	Q: Czy w zdaniu była mowa o oponach? A: Tak
\item Podczas przerw między rundami graczom źle doradzały trenery.\\
	Q: Czy w zdaniu była mowa o przeprowadzce? A: Nie
\item Na forum ekonomiczne przyjechały prezydenty.\\
	Q: Czy w zdaniu wspomniano forum ekonomiczne? A: Tak
\item Z naszej parafii na misje pojechały księdze.\\
	Q: Czy w zdaniu była mowa o chórze? A: Nie
\item Największy stół położony przy oknach zajęły proboszcze.\\
	Q: Czy stół znajdował się przy drzwiach? A: Nie
\item Z nożycami do rozcinania auta przyjechały strażaki.\\
	Q: Czy w zdaniu była mowa o rozcinaniu samochodu? A: Tak
\item Na salę operacyjną weszły chirurgi.\\
	Q: Czy w zdaniu wspomniano o pielęgniarkach? A: Nie
\item Na samym końcu do samolotu weszły piloty.\\
	Q: Czy w zdaniu była mowa o stewardessach? A: Nie
\item Zazwyczaj z potłuczonymi nosami do nas przychodziły boksery.\\
	Q: Czy w zdaniu wspomniano o urazie kręgosłupa? A: Nie
\item Większe dotacje na organizację festiwalu wywalczyły dyrektory.\\
	Q: Czy w zdaniu była mowa o organizacji festiwalu? A: Tak
\item W trójkąt stojący sto metrów za naszym samochodem wjechały policjanty.\\
	Q: Czy w zdaniu była mowa o straży pożarnej? A: Nie
\item Przez hałas na ulicy z zakładu wyszły rzeźniki.\\
	Q: Czy w zdaniu wspomniano o rzeźnikach? A: Tak
\item Przed wejściem do zbiornika wodnego nas zatrzymały strażniki.\\
	Q: Czy w zdaniu wspomniano o zbiorniku wodnym? A: Tak
\item Proszku na szczury w piwnicy podsypały dozorce.\\
	Q: Czy w zdaniu wspomniano o parterze budynku? A: Nie
\item Klimatyzowanymi limuzynami jeździły prezesy.\\
	Q: Czy w zdaniu była mowa o limuzynach? A: Tak
\item Na nagrania do Ameryki Południowej poleciały reżysery.\\
	Q: Czy w zdaniu wspomniano o Azji? A: Nie
\item Ze sprzedaży swoich produktów na miejskich targowiskach żyły rolniki.\\
	Q: Czy w zdaniu była mowa o targowiskach miejskich? A: Tak
\item Blokujący wodę wał postawiły inżyniery.\\
	Q: Czy w zdaniu wspomniano o moście? A: Nie
\item W tym budynku pracowały szewce.\\
	Q: Czy w zdaniu była mowa o szewcach? A: Tak
\item Największy problem ze znalezieniem mojego domu miały kuriery.\\
	Q: Czy zdanie opisywało zwrot nadanej przesyłki? A: Nie
\item Dużo kłopotu z naszymi bagażami miały celniki.\\
	Q: Czy w zdaniu była mowa o bagażach? A: Tak
\end{enumerate}

\subsection*{Control items:}
\begin{enumerate}
\item \textbf{context} Nie wszystkie dzieci chciałyby mieć wieczne wakacje.\\
  \textbf{correct} Syn sąsiadki jest szczęśliwy że idzie do szkoły.\\
  \textbf{incorrect} Syn sąsiadki jest szczęśliwa że idzie do szkoły.\\
	Q: Czy syn sąsiadki chodzi do szkoły? A: Tak
\item Moja córka wychowuje dzieci sama.\\
	Mąż mojej córki jest szczęśliwy w Hollywood.\\
	Q: Czy w zdaniu wspomniano Hollywood? A: Tak
\item Kiedy wróciłem z pracy, na podjeździe stało jego auto.\\
	Kochanek mojej żony został przyłapany.\\
	Q: Czy według zdania żona miała kochanka? A: Tak
\item Na opiekunkę możemy wziąć kogoś z rodziny.\\
	Kuzynka mojego ojca jest miła dla dzieci.\\
	Q: Czy według zdania ojciec ma kuzynkę? A: Tak
\item Widzę coraz więcej starszych mężczyzn spotykających się z dwudziestolatkami.\\
	Partnerka mojego wujka jest młoda.\\
	Q: Czy wspomniana w zdaniu kobieta jest młoda? A: Tak
\item Po drodze do Hiszpanii możemy się zatrzymywać u znajomych.\\
	Córka mojego przyjaciela wyjechała do Niemiec.\\
	Q: Czy w zdaniu wspomniano o Niemczech? A: Tak
\item Radio w samochodzie mi nie działa.\\
	Czubek anteny zahaczył o sufit w tunelu.\\
	Q: Czy według zdania samochód jechał przez most? A: Nie
\item Jeszcze nie skończyłem sprzątać.\\
	Fragment posadzki jest poplamiony.\\
	Q: Czy posadzka jest czysta? A: Nie
\item Źle mi się spało dzisiaj w nocy.\\
	Blask latarni mnie budził.\\
	Q: Czy powodem budzenia był hałas u sąsiadów? A: Nie
\item Usłyszałam hałas w kuchni.\\
	Pokrywka garnka spadła na ziemię.\\
	Q: Czy słoik upadł na ziemię? A: Nie
\item Spodobała mi się taka jedna piosenka, ale nie mogę sobie przypomnieć tytułu ani nazwy zespołu.\\
	Okładka albumu miała zdjęcie kwiatów.\\
	Q: Czy według zdania na okładce była twarz? A: Nie
\item To miejsce nie zachęca klientów do wejścia do środka.\\
	Witryna sklepu wygląda na starą.\\
	Q: Czy zdanie wspomina o szyldzie sklepu? A: Nie
\end{enumerate}

\subsection*{Filler items:}
\begin{enumerate}
\item \textbf{context} Bardzo lubię zajmować się wystrojem wnętrz.\\
	\textbf{item} Mój brat kupił mieszkanie, które kazałem mu wyremontować.\\
	Q: Czy brat kupił dom? A: Nie
\item Pracownicy sklepu niezbyt mi pomogli z usterką.\\
	Otrzymałem instrukcję obsługi, którą powiedziano mi żebym przeczytał.\\
	Q: Czy w zdaniu mowa o wierszu? A: Nie
\item Jutro przyjeżdża do mnie rodzina i chcę ich dobrze ugościć.\\
	Na obiad zaplanowałem danie rybne, które potrafię przygotować.\\
	Q: Czy danie rybne ma być na kolację? A: Nie
\item Po tym kursie bardzo dobrze znam język grecki.\\
	Przeczytałem grecką powieść i zastanawiam się, jak ją przetłumaczyć.\\
	Q: Czy w zdaniu wspomniano o niemieckiej powieści? A: Nie
\item Szukam agencji opiekunek do dzieci.\\
	Mam córkę i zastanawiam się, gdzie znaleźć dla niej opiekę w środy.\\
	Q: Czy córka potrzebuje opieki w piątki? A: Nie
\item Muszę kupić otwieracz do wina.\\
	Kupiłem butelkę i nie wiem, jak ją otworzyć.\\
	Q: Czy w zdaniu mowa o otwieraniu drzwi? A: Nie
\item Jutro chcę zadzwonić do brata.\\
	Mam problem, którego nie potrafię rozwiązać samodzielnie.\\
	Q: Czy osoba w zdaniu może rozwiązać problem samodzielnie? A: Nie
\item Strona internetowa sklepu jest trudna do nawigowania.\\
	Potrzebuję odpowiedzi i chciałbym ją znaleźć na tej stronie.\\
	Q: Czy osoba w zdaniu zna odpowiedź? A: Nie
\item Szukam zeszytu w którym babcia zapisywała hasła.\\
	Znalazłem na strychu sejf i nie wiem jak go otworzyć.\\
	Q: Czy według zdania na strychu była skrzynia? A: Nie
\item Potrzebuję jakiejś okazji do świętowania.\\
	Kupiłem drogie wino, którego nie wiem kiedy spróbować.\\
	Q: Czy wino było tanie? A: Nie
\item Opieka nad psem nie byłaby zbyt dużą odpowiedzialnością dla mnie.\\
	Mam w domu dużo roślin, które umiem pielęgnować.\\
	Q: Czy w domu jest mało roślin? A: Nie
\item Jak będę w centrum miasta to wejdę do apteki.\\
	Przepisano mi nowe tabletki, które nie wiem gdzie kupić.\\
	Q: Czy przepisano nową maść? A: Nie
\item W weekend planuję zrobić przemeblowanie.\\
	Mam nowe biurko i szukam miejsca, gdzie je postawię w swoim domu.\\
	Q: Czy w zdaniu wspomniano o biurku? A: Tak
\item Muszę zadzwonić do sklepu.\\
	Dostałem regał na książki, ale nie wiem jak go złożyć w domu.\\
	Q: Czy w zdaniu wspomniano o regale na książki? A: Tak
\item Ta strona jest zbyt skomplikowana jak dla mnie.\\
	Mam pytanie i nie wiem, gdzie je umieścić na tym forum.\\
	Q: Czy osoba w zdaniu ma pytanie? A: Tak
\item Bardzo mi się podobało na zajęciach z ogrodnictwa.\\
	Wymieniłem się nasionami i planuję je zasadzić w ogrodzie.\\
	Q: Czy osoba w zdaniu chce zasadzić nasiona w ogrodzie? A: Tak
\item Myślę, że możemy zakończyć na dziś.\\
	Pozostała jedna kwestia i planuję poruszyć ją następnym razem.\\
	Q: Czy pozostała kwestia będzie poruszona następnym razem? A: Tak
\item Zawsze wolałem psy, ale ostatnio się przekonuję do kotów.\\
	Mój sąsiad powierzył mi kotka i powiedział, jak się nim opiekować.\\
	Q: Czy kot był powierzony przez sąsiada? A: Tak
\item Rodzice pozwolili mi wracać samemu do domu ze szkoły w pierwszej klasie.\\
	Powierzono mi klucze i pamiętam, że odłożyłem je w bezpieczne miejsce.\\
	Q: Czy klucze są w bezpiecznym miejscu? A: Tak
\item Chciałbym mieć trochę spokoju w domu w sierpniu.\\
	Mam syna i rozważam wysłanie go na kolonie.\\
	Q: Czy w zdaniu wspomniany jest syn? A: Tak
\item Ostatnio mam więcej czasu na moje hobby.\\
	Skończyłem kolejną rzeźbę i zastanawiam się gdzie ją postawić.\\
	Q: Czy w zdaniu wspomniano o rzeźbie? A: Tak
\item Zanim zaczniemy gotować, pojadę na szybkie zakupy.\\
	Brakuje mi jednego składnika i wiem gdzie go kupić.\\
	Q: Czy osoba w zdaniu wie gdzie kupić brakujący składnik? A: Tak
\item Chcę spać latem na działce, ale jest pełno komarów.\\
	Kupiłem siatkę na komary, którą wiem jak zamontować.\\
	Q: Czy w zdaniu mowa o siatce na komary? A: Tak
\item Czekam na jakieś zniżki na loty do Włoch.\\
	Mam kilka dni urlopu i zastanawiam się kiedy je wykorzystać.\\
	Q: Czy w zdaniu wspomniano o urlopie? A: Tak
\end{enumerate}


%%% Local Variables:
%%% mode: LaTeX
%%% TeX-master: "../main"
%%% End:
